% 第一章第四节：明确事务范围、金额单位与推演条件。
前面这套处理过程依赖几个具体条件：四张表位于同一数据库，一笔购买使用站内余额付款，金额以分保存。接下来逐项检查这些条件。改变数据部署、付款路径或金额单位，都可能影响刚才的推导。

\subsection*{一笔事务覆盖哪些数据}

本章四张表位于同一个 PostgreSQL 数据库中，由同一连接的一笔事务写入。应用可以使用多个线程，也可以部署多个进程。只要这些执行者访问同一份权威数据，并遵守条件扣款与整体提交的规则，购买仍按前面说明的事务边界处理。

如果把账户和订单放进两个独立数据库，并分别提交，一边成功、另一边失败就会留下部分结果。此时已经出现两个提交边界，需要另行处理它们之间的故障。本章将四张表保存在同一数据库中，因此可以用本地事务整体提交或撤回。

详情查询也从主数据库读取，与购买写入使用同一份权威记录。若改从异步只读副本查询，复制延迟可能使刚提交的购买暂时不可见。采用这条读取路径时，需要把复制延迟纳入接口与查询约定。

数据库持久性同样有前提。我们使用数据库正常提供的事务与恢复能力，按配置保留已确认的提交。进程退出、数据库正常重启和人为把数据库恢复到昨天的备份，是不同操作。恢复旧备份会让今天的购买记录消失，后面讨论恢复时再检查这种历史变化。

\subsection*{一笔购买只走一条付款路径}

Alice 的 3000 分由站内余额支付。购买成功需要余额修改、订单和扣款流水共同提交。外部支付和后续履约各有自己的处理过程，接入时需要分别确定结果如何确认、如何保存。

这个范围让回滚容易解释：尚未提交的变化都在同一数据库事务内。外部平台的账本不会因为本地 \code{ROLLBACK} 撤回。把一个支付 HTTP 调用放在 \code{BEGIN} 和 \code{COMMIT} 之间，也无法把平台变成本地事务的一部分。

图~\ref{fig:ch01-boundary} 标出了本地事务包含的数据。接入外部支付后，付款动作由外部平台处理，本地记录平台返回的结果与编号。这两边各自提交，就会出现一边已经完成、另一边还在等待的时刻，后续章节需要逐一处理这些情况。

\begin{figure}[htbp]
\centering
\begin{tikzpicture}[x=1cm,y=1cm]
\node[booknode,text width=2.5cm] (app) at (0,0) {应用线程\\或多个进程};
\node[booknode,text width=2.2cm] (a) at (5,1) {账户};
\node[booknode,text width=2.2cm] (p) at (8.3,1) {购买};
\node[booknode,text width=2.2cm] (o) at (5,-1) {订单};
\node[booknode,text width=2.2cm] (d) at (8.3,-1) {扣款流水};
\begin{scope}[on background layer]
\node[draw=Ink!55,fill=Paper!35,rounded corners=5pt,fit=(a)(p)(o)(d),inner sep=13pt] (tx) {};
\end{scope}
\node[booknote] at (6.65,2.15) {同一数据库、同一连接的一笔事务};
\draw[bookarrow] (app.east)--node[above,font=\scriptsize]{业务写入}(tx.west);
\node[bookresult,text width=2.4cm] (ext) at (12.6,0) {外部支付平台\\使用独立账本};
\draw[bookalert,dashed] (tx.east)--node[above,font=\scriptsize]{独立调用}(ext.west);
\node[booknote,text width=3.4cm] at (12.4,-2) {本地回滚的作用\\止于数据库事务边界};
\end{tikzpicture}
\caption{事务的范围由参与写入的数据与连接决定。增加应用进程仍可共用本地事务规则，外部支付则引入新的边界。}
\label{fig:ch01-boundary}
\end{figure}


\subsection*{从显示金额到整数金额}

金额字段统一使用分。30.00 元进入服务端目录时记录为整数 3000，SQL 比较 10000 是否至少为 3000，更新时执行整数减法，响应也返回 3000。页面展示时再格式化成 30.00 元。

这条转换链应当采用同一约定。接口和数据库中的 \code{amount\_minor} 都表示分，页面再按分与元的换算关系显示。若一处把 30 解释为元，另一处把它解释为分，就会相差一百倍；字段单位应在接口定义中明确说明。

可以直接用整数的商与余数格式化：3000 除以 100，商为 30、余数为 0，显示为 \code{30.00}；3050 得到 30 和 50，显示为 \code{30.50}。金额进入余额运算前不需要先变成二进制浮点数，再乘 100 猜测怎样取整。

整数也需要范围检查。目录金额必须为正，不能超过业务允许的上限，并且能够放进 \code{bigint}。将来支持数量、税费、折扣时，还要说明在哪里取整、余数怎样分配。本章只有一件商品数量一，暂时没有这些运算。

币种固定为 CNY，使 3000 分有唯一解释。若同时支持人民币和其他货币，就需要在账户、目录、订单和流水之间核对币种，不能拿一种币种的余额直接支付另一种币种的价格。固定币种让我们先把单币种扣款的每一步检查清楚。

\subsection*{顺着记录检查一轮调用}

把前三节的调用按顺序放在一起，可以检查每次处理后留下的状态。账户 1 完成一次 3000 分购买和一次余额不足的 8000 分购买后，共有两条购买记录、一张订单和一条流水，余额为 7000 分。随后再完成一次 3000 分购买，余额降为 4000 分。

下表列出这三个步骤的累计结果。余额单位为分，最后一列分别列出订单数与流水数。逐行查看时，可以同时核对每次扣款和新增记录之间的对应关系。

\begin{center}
\begin{minipage}{\linewidth}
\centering\small\captionsetup{hypcap=false}
\captionof{table}{三次购买后的累计状态}
\begin{tabularx}{\linewidth}{Xcccc}
\toprule
\textbf{步骤} & \textbf{HTTP} & \textbf{余额} & \textbf{累计购买} & \textbf{订单/流水} \\
\midrule
买 3000 分相纸 & 201 & 7000 & 1 & 1/1 \\
买 8000 分相纸 & 409 & 7000 & 2 & 1/1 \\
再买 3000 分相纸 & 201 & 4000 & 3 & 2/2 \\
\bottomrule
\end{tabularx}
\end{minipage}
\end{center}

两次成功购买应生成不同的购买、订单和扣款编号。同一账户能够查询自己的购买，其他账户查询时返回 404。若在写流水时发生 SQL 错误，当前事务应回滚；对于从 10000 分和空历史开始的这笔购买，回滚后的余额恢复为 10000 分，购买、订单和流水均为空。

还可以按购买编号关联三张业务表，一次检查购买金额、订单金额、扣款金额和当次扣后余额：

\begin{BookCode}
SELECT p.purchase_id,
       p.amount_minor AS purchase_amount,
       o.order_id, o.amount_minor AS order_amount,
       d.debit_id, d.amount_minor AS debit_amount,
       d.balance_after_minor
FROM purchases AS p
LEFT JOIN orders AS o ON o.purchase_id = p.purchase_id
LEFT JOIN debits AS d ON d.order_id = o.order_id
WHERE p.purchase_id = :purchase_id;
\end{BookCode}

成功记录的三个金额都应为 3000，订单号和扣款号都存在。余额不足记录则只有购买金额，后面的订单与扣款字段为空。然后再执行第二节的余额合计查询，把账户当前余额和全部扣款连起来核对。

记录数量与金额要一起看。只看到一张成功订单，还无法排除错误多扣了一次；只看到余额是 7000，也无法确认扣款是否已经留下流水。把编号、金额、状态和数量接起来检查，才能解释这次购买。

重新演算时，需要同时确定初始余额和已有流水。若把当前余额单独改回 10000 分，却保留原来的 3000 分扣款流水，核对公式会得出预期余额 7000 分，与当前余额相差 3000 分。使用空历史的初始状态，或沿已有历史继续计算，都能得到明确的比较依据。

\subsection*{推演练习：两个独立账户，两种顺序}

账户 2 和账户 3 的初始余额各为 10000 分，购买、订单和流水均为空。用两个独立账户推演下面两组顺序，比较执行次序怎样影响购买结果。

先在账户 2 上依次处理三次新购买：

\noindent A：购买 3000 分的 \code{paper-001}。

\noindent B：购买 8000 分的 \code{paper-002}。

\noindent C：再购买 3000 分的 \code{paper-001}。

每次都生成新的购买编号，上一笔事务结束后再开始下一笔。先写下每一步的余额、购买状态、累计订单数和累计流水数，再计算最终扣款金额合计。

按 A、B、C 顺序处理，A 扣掉 3000，余额为 7000，产生第一张订单与第一条流水。B 需要 8000，条件更新返回零行，只提交余额不足的购买结果。账户保持 7000，订单和流水数量没有增加。C 扣掉 3000，余额变为 4000，产生第二张订单和第二条流水。

最终购买三条，状态依次为成功、余额不足、成功；订单两条，流水两条；扣款合计 6000，\code{10000 - 6000 = 4000}。余额不足的 B 在购买表中占一条记录，在扣款合计中占零。

再用初始状态相同的账户 3 改变顺序：A 和 B 同时开始，由 B 先完成账户更新并提交，最后执行 C。这里指定 B 先获得更新机会，用来检查这一条具体的并发执行历史。

B 先把余额从 10000 分减到 2000 分。A 的更新随后在新行版本上重新判断，2000 小于 3000，进入余额不足分支。C 再执行时也因余额不足结束。最终只有 B 的一张订单和一条 8000 分流水，账户余额为 2000 分。

两种顺序下，成功的购买不同，各自的余额变化却都与扣款流水一致。条件更新依据取得账户更新机会时的余额作出判断，事务保证相关记录共同提交。这里比较的是不同购买的执行次序；下一章将继续检查同一购买重发请求的情况。

最后比较同一初始状态下的提交与回滚。账户有 10000 分、业务记录为空时，开始一笔 3000 分购买。若在扣款和订单插入后回滚，账户恢复为 10000 分，购买、订单和流水均为空；若继续补齐流水和购买结果并提交，账户为 7000 分，三类业务记录各有一条。两种结果的差别取决于整笔事务最终是提交还是撤回。

Alice 现在可以用购买编号查到自己买的商品、成交金额、订单和扣款记录；余额不足的购买也有自己的查询结果。下一章保留同一条成功路径，只在提交后、响应发出前中断连接。数据库里的订单和流水已经保存，Alice 却没有收到它们的编号。她再次点击购买时，服务会怎样处理，就从这里继续看。

本章使用的数据库行为可参阅 PostgreSQL 18 手册中的\href{https://www.postgresql.org/docs/18/tutorial-transactions.html}{事务教程}、\href{https://www.postgresql.org/docs/18/transaction-iso.html#XACT-READ-COMMITTED}{Read Committed 隔离级别}与 \href{https://www.postgresql.org/docs/18/dml-returning.html}{RETURNING 子句}。
